\section{Consola de depuración}

Una web mínima en Next.js para probar los dos agentes en vivo, con datos reales. No es
un producto: es el instrumento con el que se verifica que un agente no está inventando.

\subsection{La regla de verificación}

Por cada pregunta, la consola dibuja la traza completa: cada turno del modelo con su
texto y las herramientas que decide llamar, cada ejecución real con su entrada, su
\textbf{salida cruda}, su tiempo en milisegundos y su error si lo hubo, y al final la
respuesta redactada con el consumo de tokens y el costo.

\begin{clave}[title=Regla de oro]
Todo número de la respuesta final tiene que aparecer en la salida de alguna herramienta.
Si aparece un número que no está en ninguna salida, el modelo lo alucinó.
\end{clave}

\subsection{Camino de un request}

\begin{center}
\small
\cd{Navegador} $\to$ \cd{/api/analizador/*} $\to$ \cd{:8010} $\to$ Supabase (RO) \\[2pt]
\cd{Navegador} $\to$ \cd{/api/pronostico/*} $\to$ \cd{:8000} $\to$ store (RO)
\end{center}

El navegador \textbf{nunca} habla directo con los servicios Python ni ve las API keys.
Todo pasa por manejadores de ruta del lado servidor, que reenvían la petición inyectando
la clave. Las claves viven sólo en la configuración del servidor.

Los endpoints que la consola necesita se agregaron a los propios agentes, no se
reimplementaron en Node. Es una decisión de fuente única de verdad: si el lazo del LLM se
reimplementara en el cliente, habría dos lazos que mantener sincronizados y la traza
dejaría de reflejar lo que realmente corre en producción.

\subsection{Control de acceso}

El \emph{middleware} decide si un request pasa; no sabe firmar cookies ni renderizar el
login. Las páginas redirigen al login; las rutas \cd{/api/*} responden 401 en JSON,
porque un \cd{fetch} no puede hacer nada útil con una redirección a HTML.

\begin{aviso}[title=Fallar cerrado]
Sin \cd{DEBUGGER\_PASSWORD} configurada, en producción la consola responde \textbf{503}
en lugar de abrirse. Fallar abierto fue exactamente lo que la dejó expuesta en Amplify
en su momento. En desarrollo sí deja pasar, para no estorbar el trabajo local.
\cd{DEBUGGER\_SESSION\_SECRET} firma la cookie: rotarlo cierra todas las sesiones sin
cambiarle la contraseña al equipo.
\end{aviso}

\subsection{Dos formas de levantarla}

\cd{./dev.sh} monta todo local: analizador, pronóstico y la web. \cd{./consola.sh} abre
un túnel SSH contra los agentes que \emph{ya corren} en la EC2 ---que siguen escuchando
sólo en el loopback del servidor, sin exponer nada--- elige puertos libres (8000, 8010 y
3000 suelen estar ocupados en una máquina de trabajo), sincroniza las URLs con los
puertos de esa corrida y cierra el túnel al salir.
